\clearpage
\fancyhead[L]{\small MERIDIAN CIRCLE.}
\hypertarget{native-434-part-lower}{}
\begin{center}
{\Large MERIDIAN CIRCLE.}

\rule{0.6\linewidth}{0.4pt}

Göttingische gelehrte Anzeigen. No. 91. pp. 905..912. 5 June 1820.

\rule{0.6\linewidth}{0.4pt}
\end{center}

The new \textsc{Reichenbach} meridian circle, which our observatory received last year, has already been set up and in use for some time. Before reporting on it publicly, however, we thought we should wait until we had gained considerable experience, so that we could judge more confidently what an instrument of a design never used before could do.

The instrument serves both as a transit telescope and as an instrument for measuring altitudes. For the former purpose it has all the features of the finest ordinary transit telescopes, so those features need no detailed description. The telescope has a focal length of 5 Paris feet and an aperture of 4 Paris inches. The four eyepieces magnify 68, 86, 120 and 170 times. Councillor \textsc{Gauss} uses the highest power almost exclusively. The reticle had to be replaced here because some of its threads had gone slack. It now consists of seven vertical and two horizontal spider threads. Stars on the equator cross the intervals between the vertical threads in $\cdot 14$ seconds. The horizontal threads are only $7''6$ apart.

The axis, 33 Paris inches long, carries two concentric circles at one end. Their faces on the side away from the telescope lie almost in one plane. The outer circle is \textit{fixed} to the axis and therefore turns with the telescope. It carries the scale, which is graduated directly at intervals of 3 \SourceLeaf{435}
minutes. The inner circle (the alidade circle) could turn about the axis were it not for a stop fitted to the pier. The stop allows only a slight, fine movement for adjusting the spirit level fixed to the circle. The four indices on the alidade circle, each $45^\circ$ from the vertical, have verniers that subdivide each main division into 90 parts, giving intervals of 2 seconds. Smaller fractions can still be estimated. At the point where they are read, the circles are 35 Paris inches in diameter. A barely perceptible gap separates the circles without their touching. The face of the graduated circle projects ever so slightly above the face of the alidade circle, which is why the reading microscopes are deliberately set at a slight angle. This instrument shares these features with others by \textsc{Reichenbach}. Three balancing arrangements, one for the instrument as a whole, one for the alidade circle and one for the telescope, are intended to cancel the pressure of the pivots on their bearing cups and of the alidade circle on the axis, as well as the flexure of the tube. A deflection of one Paris inch corresponds to $22''$ on the hanging spirit level used to level the axis, and to $17''6$ on the main spirit level. The latter serves to keep the alidade circle in a fixed position, or to measure and allow for very small shifts. (Councillor \textsc{Gauss} usually touches the adjusting screw only when the shift exceeds $2''$, which rarely happens.)

The instrument can be reversed easily and safely with a machine designed very well for the purpose. In measuring altitudes, the reversal serves first of all to determine the collimation error and to convert the measured zenith distances into absolute ones. It may seem questionable, however, to assume that the collimation error stays exactly constant over a long period and through considerable changes in temperature. Unless the instrument is reversed very often, it will therefore be better to refer the observations to the pole than to the zenith. The remarkable optical power of the telescope makes it all the easier to determine the position of the pole on the circle from circumpolar stars. In this respect it is hardly inferior to the telescope on \textsc{Reichenbach}’s transit instrument (this journal, 1819 [p. 422 of this volume]), even though the latter has a longer focal length and a somewhat larger aperture. Any small difference that remains is offset by the fact that the meridian circle can be set beforehand with the utmost precision, so that one always has to give one’s full attention to only a very small part of the field of view. Among the stars Councillor \textsc{Gauss} observed in daylight, many are of the fourth magnitude. 1 Draconis \textsc{Hevelii} and $\omega$ Cephei \textsc{Hevelii}, also observed in daylight close to noon, are only of the fifth magnitude. The design of the reticle makes daytime observations of such tiny points, if they are visible at all, as reliable as observations of the brighter stars.

\Needspace{40\baselineskip}
To show the splendid agreement this instrument can achieve, we give the results of all determinations of the position of the pole on the instrument from the opening of the observing log to the first reversal. They are drawn from all opposite culminations of circumpolar stars observed immediately one after the other. The agreement among results from one and the same star shows the precision with which the observations can be made and shows that the collimation error stayed constant during that time. The equally close agreement among results from different stars shows the excellence of \textsc{Reichenbach}’s graduation, all the more admirable given the moderate size of the circle.

\Needspace{29\baselineskip}
\SourceLeaf{436}
\begin{center}
\begin{tabular}{l|l|c|c}
1820 & Star name & Magnitude & Position of the pole \\
\hline
Feb. 26 & $\alpha$ Cephei & 3 & $321^\circ\,29'\,32''14$ \\
 & $\beta$ Cephei & 3 & $31.89$ \\
 & $\delta$ Draconis & 3 & $32.29$ \\
 & $\eta$ Cephei & 3.4 & $32.24$ \\
 & $\alpha$ Cephei & 3 & $31.16$ \\
 & $\beta$ Cephei & 3 & $31.65$ \\
27 & $\delta$ Draconis & 3 & $31.39$ \\
 & $\eta$ Cephei & 3.4 & $31.90$ \\
 & $\delta$ Draconis & 3 & $31.54$ \\
 & $\eta$ Cephei & 3.4 & $31.63$ \\
28 & $\gamma$ Cephei & 3 & $31.54$ \\
 & $\tau$ Cephei & 4 & $32.74$ \\
 & $\delta$ Cephei & 4 & $31.12$ \\
Mar. 9 & $\delta$ Draconis & 3 & $31.18$ \\
 & $\beta$ Cephei & 3 & $32.18$ \\
 & $\tau$ Cephei & 4 & $31.42$ \\
 & $\iota$ Cephei & 4 & $31.52$ \\
10 & $\beta$ Cephei & 3 & $32.05$ \\
11 & 1 Dracon. Hev. & 5 & $31.97$ \\
18 & 1 Dracon. Hev. & 5 & $32.04$ \\
 & 3 Lacertae & 4 & $32.32$ \\
 & $\alpha$ Ursae maj. & 2 & $31.62$ \\
19 & $\alpha$ Ursae min. & 2 & $32.81$ \\
\end{tabular}
\end{center}

\par\indent The mean is $321^\circ\,29'\,31''84$. Not one of these 23 determinations differs from it by a full second.

\par\indent Insofar as observations with this instrument are referred to the pole instead of to the zenith, the most precise knowledge of the altitude of the pole is really of secondary importance. Where a single second or even fractions of a second matter, the greatest care is needed, and one cannot expect the collimation error to remain \textit{as constant} indefinitely as the table above shows. The most precise determination of the altitude of the pole is therefore best made by combining only observations that immediately preceded and immediately followed a reversal. Councillor \textsc{Gauss} therefore does not venture to use his observations so far for \textit{this} determination. After the first reversal the weather turned unfavorable for quite some time, and at the second, on 13 April, a chance slight disturbance of the reticle made the \textit{complete} equality of the collimation error somewhat doubtful. The results from those observations nevertheless agree very closely with one another and with a result obtained by another method, which we shall give shortly. Combining the position of the pole in the first period with that in the second ($38^\circ\,25'\,53''\!54$) gives an altitude of the pole of $51^\circ\,31'\,49''\!15$. Combining 13 observations of circumpolar stars immediately before the second reversal with the opposite culminations of the same stars observed immediately afterward gives $51^\circ\,31'\,47''\!86$.

It is well known that a circumstance has drawn considerable attention from astronomers in recent years. Although only a couple of seconds are involved, it is of the greatest importance both for the art of observation and for a wide range of astronomical results. We mean the small differences that have appeared in determinations of stellar declinations, the obliquity of the ecliptic and the altitude of the pole when different, though otherwise excellent, instruments are used. There is no doubt that these differences come from the \emph{action of gravity} on the various parts of each instrument, even though up to \SourceLeaf{437}
now it has been possible neither to establish fully and with certainty how it acts nor to judge definitely which instrument gave the correct result and which gave the incorrect one. We still know very little about how much metals bend, and it seems to us too bold to deny, for any instrument however built, the possibility that gravity might have an appreciable effect on its parts and thereby on the observations, without sufficient independent evidence for such a claim. In the meridian circle here, the great instrument maker has indeed done all he could to eliminate the flexure of the telescope by ingenious balancing. It might still be doubted, however, whether this eliminated all flexure or made its effects completely imperceptible. The only direct way to test this seems to be to combine direct observations of a celestial body with observations of its image seen in an artificial horizon. Such observations of course have to be repeated often to clear up this delicate matter. Councillor \textsc{Gauss} has already begun to make such observations by observing the Pole Star in a reflecting surface of water. The most striking evidence of the astonishing optical power of the telescope is probably that the upper culmination, which occurs in daylight, could also be observed quite well in this way. The result of the first \textit{complete} observation of this kind was as follows:

\Needspace{22\baselineskip}
{\rightskip=0pt plus 3em \textit{Zenith distance of the Pole Star} (free from refraction, but including the culmination error)\par}

\begin{center}
1820 May 13.
\end{center}

\begin{center}
\begin{tabular}{l c l}
Lower culmination, direct & . . . . . . . . . . & $319^\circ\,50'\,20''\!73$ \\
Lower culmination, reflected & . . . . . . . . . & $220\ \ 5\ \ 3.94$ \\
Upper culmination, direct & . . . . . . . . . . & $323\ \ 8\ 41.51$ \\
Upper culmination, reflected & . . . . . . . . . & $216\ 46\ 44.31$ \\
\end{tabular}
\end{center}

\begin{center}
This gives the \textit{true} zenith distance
\end{center}

\begin{center}
\begin{tabular}{l c l}
at lower culmination & . . . . . . . . . . & $40\ \ 7\ 21.60$ \\
at upper culmination & . . . . . . . . . . & $36\ 49\ \ 1.40$ \\
\end{tabular}
\end{center}

\noindent and from this, since the change in declination over 12 hours is $-0''10$, the altitude of the pole is $51^\circ\,31'\,48''\!45$ at the location of the water vessel, or $51^\circ\,31'\,48''\!40$ at the center of the circle.

\par\indent Since this determination falls almost midway between the two given above, it already becomes very probable that the effect of gravity on observations with this instrument must either be entirely imperceptible or be extremely small. The declination of the Pole Star from the observations above is $0''42$ smaller than that from \textsc{Bessel}’s tables (all the observations made so far by Councillor \textsc{Gauss} give this correction $=
-0''67$).

\Needspace{17\baselineskip}
To show what the instrument can do as a transit telescope as well, we give a few observations of Mars made by Councillor \textsc{Gauss} around the time of that planet’s quadrature.

\begin{center}
\begin{tabular}{l|c|c|c}
1820 & Mean time & \shortstack{R.A. of\\1st limb} & \shortstack{Decl. of\\center} \\
\hline
Mar. 29 & $7^{\mathrm u}\,10^{\mathrm m}\,16^{\mathrm s}3$ & $114^\circ\,37'\,49''3$ & $24^\circ\,16'\,47''3$ N. \\
30 & $7\,7\,40.7$ & $114\,57\,57.7$ & $24\,12\,7.1$ \\
31 & $7\,5\,6.8$ & $115\,18\,30.1$ & $24\,7\,18.1$ \\
Apr. 5 & $6\,52\,37.3$ & $117\,6\,19.9$ & $23\,42\,5.4$ \\
10 & $6\,40\,40.4$ & $119\,2\,17.5$ & $23\,14\,30.6$ \\
11 & $6\,38\,20\,3$ & $119\,26\,18.3$ & $23\,8\,41.0$ \\
12 & $6\,36\,1.6$ & $119\,50\,41.5$ & $23\,2\,48.0$ \\
13 & $6\,33\,43.7$ & $120\,15\,15.0$ & $22\,56\,44.8$ \\
26 & $6\,5\,16.8$ & $125\,56\,10.0$ & $21\,29\,12.8$ \\
\end{tabular}
\end{center}

\Needspace{18\baselineskip}
\SourceLeaf{438}
Mr. von \textsc{Staudt} compared these observations with Mr. von \textsc{Lindenau}’s tables for Mars. He is now studying mathematics and astronomy here with outstanding zeal and success. The comparison yielded the following

\begin{center}
Differences

\begin{tabular}{ll|r|r}
 & & Right asc. & Decl. \\
\hline
Mar. & 29 & $-0''7$ & $+5''1$ \\
 & 30 & $-0.5$ & $+3.9$ \\
 & 31 & $-2.1$ & $+6.1$ \\
Apr. & 5 & $-1.3$ & $+3.9$ \\
 & 10 & $-4.9$ & $+2.9$ \\
 & 11 & $-1.9$ & $+3.6$ \\
 & 12 & $-4.8$ & $+1.9$ \\
 & 13 & $-1.8$ & $+4.2$ \\
 & 26 & $-1.5$ & $+1.5$ \\
\end{tabular}
\end{center}
